\chapter{Финальное заключение Тома X и программа Тома XI}
\label{ch:v10-final-conclusion-tome11-program}

\section{Какие задачи ставились}

Том X был открыт после масштабного no-go Тома IX. Вместо поиска ещё одной
статической энергии он ввёл геометрическую историю
\begin{equation}
 \zeta=\frac13\log\frac{N}{N_0}
\end{equation}
и поставил вопрос, может ли абсолютный физический масштаб возникнуть из
нормированного рождения ячеек, квантового хода и неравновесной геометрии.

Собственное введение Тома X сформулировало шесть операционных требований:
\begin{enumerate}
\item общий носитель истории роста, локального спектрального атласа и полей;
\item физическое происхождение нормированной меры рождения ячеек;
\item ненулевой квантовый ход или следовая аномалия относительно $\zeta$;
\item физическая связь темпа роста с космологической кривизной;
\item типизированная карта геометрического хода в локальный энергетический сектор;
\item хотя бы одно безусловное безразмерное следствие.
\end{enumerate}

Одновременно Том IX передал более сильный quantum-RG контракт: вывести
RG-инвариантный абсолютный масштаб трансмутации, встроить его в KMS/Gaussian
parent и получить физическую меру logdet. Эти два контракта нельзя смешивать:
операционный контракт проверяет новую геометрию роста, а межтомовый контракт
проверяет, решена ли исходная проблема абсолютного масштаба.

\section{Ответ операционного контракта}

Для всех шести пунктов построены точные условные архитектуры:
\begin{equation}
 \mathbf n_{X,\rm op}^{\rm cond}=(1,1,1,1,1,1)^T,
 \qquad N_{X,\rm op}^{\rm cond}=6/6.
\end{equation}
Строгий физический реестр равен
\begin{equation}
 \mathbf n_{X,\rm op}^{\rm phys}=(1,0,1,0,1,1)^T,
 \qquad N_{X,\rm op}^{\rm phys}=4/6.
\end{equation}
Дефицит
\begin{equation}
 \mathbf d_{X,\rm op}=(0,1,0,1,0,0)^T
\end{equation}
имеет rank $2$ и состоит из происхождения меры рождения и физического
моста рост--кривизна.

\section{Что удалось получить}

Устойчивое ядро Тома X состоит из следующих результатов.

Во-первых, построен общий типизированный carrier геометрической истории,
локального спектрального атласа, KMS-сектора и полей. Общий статический
масштаб отделён от внутреннего безразмерного разрешения.

Во-вторых, взаимно обратный спектральный закон
\begin{equation}
 \Sigma_{\rm in}=e^\zeta,
 \qquad
 \Sigma_{\rm out}=e^{-\zeta}
\end{equation}
получен как нулевая траектория положительного parent; ориентированный
следовой свидетель ненулевой. Тем самым безразмерный квантовый ход
относительно $\zeta$ действительно построен.

В-третьих, нормированная мера рождения, часы, четырёхмерный объём,
хопфовский KMS-цикл и неравновесный двухрезервуарный ток сведены в одну
реляционную архитектуру. Получены слепые отношения
\begin{equation}
 \frac{J_{\rm edge}}{\Omega}=\Delta\zeta,
 \qquad
 \frac{\sigma}{\Omega}=3\Delta\zeta\log2,
\end{equation}
не зависящие от общей часовой нормировки.

В-четвёртых, индуцированная ньютоновская площадь и клеточная площадь дают
безмасштабное отношение, а дискретное разрешение относительно первичной
ячейки остаётся RG-инвариантным и не требует знания абсолютного размера
ячейки.

В-пятых, particle-wrinkle ветвь доведена от Callias-профиля через
Pati--Salam $\Sigma$, hypercharge splitting и relative-Hodge construction
до полного Mathai--Quillen-аудита. Условный Thom-квартет непротиворечив,
но его физический parent закрыт строгим no-go: required shift rank равен
$8$, quadratic flat rank равен $2$, ordinary gauge translational rank
равен $0$, а новая Stückelberg-копия не унаследована.

\section{Чего получить не удалось}

Нормировка вероятностей не выбирает сами веса рождения и не задаёт
физическую секунду. Орбита
\begin{equation}
 (E,t)\longmapsto(cE,t/c)
\end{equation}
сохраняется во всех построенных clock/KMS реализациях.

Кривизна, определённая через тот же темп роста, не является независимым
якорем этого темпа. Эйнштейновский мост требует независимо полученных
$G$, температуры/энергии резервуара, объёма ячейки и stress-energy map;
физический пакет этих входов не найден.

Размерная трансмутация построена только относительно опорного масштаба.
Асимптотически свободный носитель, граничная связь и типизированный портал
не происходят из текущего родителя $K_{43}$ без нового сектора. Поэтому
абсолютный масштаб не выведен.

Не получена и физическая мера logdet. Рассмотренные фермионные и
BRST-завершения условно воспроизводят нужный determinant. Подходы
Hubbard--Stratonovich, superconnection и Mathai--Quillen требуют нового
вспомогательного носителя, сдвиговой симметрии либо семени меры.

\section{Итог более сильного quantum-RG обещания}

Для шестипунктового контракта, переданного Томом IX, строгий реестр равен
\begin{equation}
 \mathbf n_{X,\rm RG}^{\rm phys}=(1,1,0,0,0,1)^T,
 \qquad N_{X,\rm RG}^{\rm phys}=3/6.
\end{equation}
Открыты три пункта: абсолютная transmutation scale, её физическое KMS-
вложение и parent-origin logdet measure. Четыре независимо сформулированных
пакета повторного открытия дали
\begin{equation}
 \mathbf a_{\rm reopen}^{\rm phys}=(0,0,0,0)^T.
\end{equation}

\begin{theorem}[Финальный ответ Тома X]
Том X построил условно полную геометрию роста с операционным счётом $6/6$
и строгим счётом $4/6$. Он вывел общий carrier, ненулевой безразмерный ход,
типизированные реляционные карты и слепые безразмерные следствия. Однако
физическое происхождение меры рождения и независимого моста к кривизне не
получено. Более сильный quantum-RG контракт закрыт лишь на $3/6$;
абсолютный масштаб и физический logdet parent не выведены.
\end{theorem}

\begin{nogocriterion}[Финальная заморозка Тома X]
Нельзя повышать нормированную меру до механизма выбора её весов, читать
$\Lambda_{\rm growth}=3(H/c)^2$ как независимый якорь $H$, использовать
reference scale dimensional transmutation как абсолютную энергию или
считать условный Stückelberg/Mathai--Quillen carrier унаследованным.
Продолжение перебора эквивалентных scale breakers и auxiliary completions
внутри Тома X запрещено.
\end{nogocriterion}

\section{Программа Тома XI}

Следующий том не должен снова искать абсолютное число внутри той же
масштабной орбиты. Его предмет --- реляционные наблюдаемые, возмущения и
проверяемые межсекторные следствия. Программа состоит из шести пунктов:
\begin{enumerate}
\item построить общую алгебру реляционных наблюдаемых роста, причинного
ядра, KMS-потока и particle-сектора;
\item явно факторизовать общую масштабную орбиту и определить полный набор
безразмерных координат физического quotient;
\item построить линейную теорию возмущений и отклика около неравновесной
растущей геометрии;
\item получить типизированные межсекторные морфизмы между particle gaps,
энтропийным потоком и космологическими реляционными наблюдаемыми;
\item вывести scheme- и gauge-независимые blind relations вместе с
устойчивостью и областью применимости;
\item заморозить предсказания до сравнения с внешними данными и провести
явный протокол фальсификации.
\end{enumerate}

Матрица зависимостей этой программы треугольна, имеет rank $6$ и
determinant $1$. Спецификация равна $6/6$, construction status при открытии
равен $0/6$. Том XI допускается как новая программа, а не как готовая
теория.

\begin{protocol}[Финальный статус Тома X]\begin{enumerate}
\item операционный условный контракт: \textbf{$6/6$};
\item операционный строгий контракт: \textbf{$4/6$};
\item строгий quantum-RG контракт Тома IX: \textbf{$3/6$};
\item физические reopening packages: \textbf{$0/4$};
\item абсолютный масштаб: \textbf{не выведен};
\item физический Mathai--Quillen parent: \textbf{не выведен};
\item Том X: \textbf{завершён и физически заморожен};
\item спецификация Тома XI: \textbf{$6/6$};
\item построение Тома XI: \textbf{$0/6$};
\item следующий гейт: \textbf{допуск общего carrier реляционных наблюдаемых}.
\end{enumerate}\end{protocol}

Машинный сертификат сохранён в
\path{s2t_v10_final_conclusion_and_tome11_program_gate_results.json}.